CPU emulators are widely used in security infrastructure, from malware sandboxes to cross-architecture exploit development. Despite this, their accuracy on instruction corner cases remains difficult to validate because many subtle CPU behaviors only appear under unusual fault, flag, or machine-state conditions. This second semester extends the previous x86-64 emulator study from a small set of preliminary probes to a broader comparison across targeted edge cases and instruction-fuzzing workflows.

Specifically, the project now compares native Linux execution against Blink, QEMU TCG, and Box64 using process-backed probes, while also evaluating Icicle, Unicorn, MWEMU, and KUBERA through shellcode-oriented runners. The new tests cover x87 status-word behavior, flag materialization, timestamp behavior, signed division faults, SSE alignment faults, MXCSR rounding, and descriptor-table or machine-status instructions. Moreover, the project now includes \texttt{sandsifter\_rs}, a Rust sandsifter-style fuzzer used to search for unusual instruction byte sequences, backend mismatches, rare instructions, and potential hidden-instruction candidates.

The current results show that every emulator differs from native hardware in at least one meaningful way. Some differences are strong correctness bugs, such as missing \texttt{IDIV} overflow faults or incorrect MXCSR rounding, while others are strong emulator fingerprints, such as synthetic system-state values. Consequently, the project now provides both deterministic regression-style probes and a broader discovery tool for future emulator accuracy research.
